\section{The Bond and the Auction in Detail}\label{app:bond}

This appendix gives the rules behind the bond of \cref{sec:bond} and the auction of \cref{sec:auction}.

\subsection{The verification game and solver pay}

\paragraph{Proving a failed opening.}
Given a key $K$ with $\Hcommit(\ctx, K) = C$, a contract checks the envelope's authentication tag directly, which needs no decryption and costs work linear in the envelope's length, which admission caps (\cref{app:auction}). A failing tag establishes malformation at once. Chains that do not lead to the committed key are harder: the evidence is the whole computation of $\Open$, $N + O(n)$ steps, which no contract can repeat. An optimistic verification game, as in TrueBit~\cite{TR17} and Arbitrum~\cite{Arbitrum18}, can establish it in about $\log_2 N$ rounds. We outline one here; a deployment needs a full specification and its gas costs (\cref{sec:open}).

\paragraph{The game.}
A \emph{state} of $\Open$ is the solver's checkpoint (\texttt{SPEC.md}, Section~6.4: the chain index, the steps done in that chain, the chain state and the accumulator) together with the current link, a few hundred bytes. The contract computes the hash of the initial state from the puzzle itself. A claimant stakes and posts the final state, whose key $K'$ the contract checks to satisfy $\Hcommit(\ctx, K') \neq C$. Within a challenge period, anyone may stake and dispute the claim. Each challenger then plays its own game against the claimant: the claimant posts the hash of the state at the midpoint of the interval still in dispute, and the challenger says whether it agrees, which halves the interval.

When one step remains, the claimant opens the state both sides agreed on, which for the first step is the initial state. The contract recomputes its successor, one SHA-256 compression or, at a chain boundary, a few labelled hashes, and compares the result's hash with the claimant's hash for the end of the interval: the claimant wins that game if they match, and the challenger otherwise. A party that misses a move's deadline loses the game, and the loser forfeits its stake. For the reference puzzle, $N + O(n)$ is just above $2^{33}$, so a game takes 34 rounds of two moves, each a transaction carrying at most one 32-byte hash, and one final step. Naraghi assumes a similar mechanism, without analysis, for the pads of blind simultaneous release~\cite{Naraghi26}.

\paragraph{Who decides.}
A claim stands only if no one challenges it within the period, or if the claimant wins every game. Games run in parallel, so a challenger colluding with the claimant that loses on purpose cannot defeat an honest one. Since $\Open$ is deterministic and its input public, anyone who has solved a well-formed puzzle can challenge a false claim, so a sender that went offline needs no one to act for it. This is the usual assumption of optimistic verification: at least one honest party solves the puzzle within the challenge period and gets each move included before its deadline. A key that matches $C$ does not refute a claim, since a malicious sender knows such a key; posted before the claim stands, it opens the transaction, ends the games, returns every stake and returns the bond.

\paragraph{Paying solvers.}
A sender that withholds its key can still publish it after the deadline, so a reward for the first matching key would go to the sender and leave the solvers that did the work unpaid, and a solver of a malformed puzzle has no key to post.

The protocol therefore assigns each missed reveal to a solver drawn from a staked registry. Within a deadline $D_s$, the solver must submit a key matching $C$, evidence that the tag fails under such a key, or a malformation claim; a key posted by anyone also counts. A timely claim keeps the assignment until the claim is decided. The solver is paid from the forfeited deposit, and recovers its stake, for a payable outcome: a valid opening, a failing tag, or its claim standing. A solver that submits nothing in time, or whose claim fails, loses its stake, and the transaction passes to the next solver. Designation prevents outsiders from copying an assignment's reward, but does not exclude solvers that the sender controls. The effective penalty is therefore $d - r$, or $d + B - r$ when malformation is established, where $r$ is every payment the sender or its coalition can recover; a forfeited bond is burned, so that $r$ is at most the solver's payment.

\paragraph{Denial of service.}
A missed reveal can induce nearly a full solve while costing only the net deposit penalty $d - r$: a sender whose chains are malformed, but who holds a key matching $C$ and a valid envelope under it, can post the key just before the solvers finish, which opens the transaction and returns the bond. The bond is additionally forfeited only when malformation is established before a successful opening. The net deposit must therefore pay for a solve at the admission limits, together with any claim, game and redundant solver. For the reference puzzle, a solve takes about 3.3 minutes of a laptop core, or 30\,s of a circuit as fast as the assumed attacker; on rented laptop-class cores at four cents per core-hour it costs about a fifth of a cent.

Deposits price demand; they do not by themselves provide solver capacity, and capacity at attacker speed is scarcer. An attacker that exhausts it delays forced openings; it does not make any transaction readable earlier, since secrecy rests on the delay, not on the solvers. A cap on open puzzles per sender does not bound the total load when the attacker uses many accounts; admission limits and deposits do.

\subsection{The auction rules}\label{app:auction}

The auction sells, in each slot, a fixed quantity of one asset for a quote asset at a uniform price. Its parameters are the deposit $d$, the bond $B$, a reveal window of $W$ slots, a solver deadline $D_s$, a challenge period $D_c$ and a move deadline $\tau$.
\begin{enumerate}[leftmargin=1.5em]
  \item \textbf{Supply and bids.} Before batch $k$ opens, the seller deposits its quantity $Q_k$ in the contract, so each batch settles from its own supply and escrows. A bid is a plaintext $(k, p, x)$: the batch, a unit price and a quantity, each a 32-byte unsigned integer, 96 bytes in all. It is \emph{valid} if its batch is $k$, $p$ is at least the reserve price, $x \le Q_k$, and $p x \le E$, where $E$ is the escrow posted in plaintext with the bid.
  \item \textbf{Admission.} The contract accepts a sealed bid for batch $k$ only in the block of slot $k$, with its escrow, $d$ and $B$, and only with the reference work and fixed limits: $nm = 2^{33}$, $n \le 16$, $b = 24$, an envelope of the bid's fixed length, and at most a fixed number of bids per batch. Checking the commitment, deriving the subkey and checking the tag then take at most 17 SHA-256 compressions.
  \item \textbf{Exposure.} A bid is exposed when it first leaves the bidder, to a builder or a mempool, and the set of bids is fixed when slot $k$'s block is final. The bidder sends its bid only to the builders of slot $k$, during that slot. The auction assumes the fast finality planned for lean consensus, under which the set is fixed within the reference $T = 30$\,s of exposure (\cref{tab:choosing-t}). Under today's finality the delay would have to cover about 13 minutes, with the larger work and costs of the last row of \cref{tab:choosing-t} (\cref{sec:reorgs}).
  \item \textbf{Reveal.} The reveal phase begins once slot $k$'s block is final, so no key is published before the set of bids is fixed. Within the next $W$ slots, the bidder publishes its key. The contract checks it against $C$, checks the tag, decrypts the bid and returns $d$ and $B$. An invalid bid is excluded and forfeits $d$.
  \item \textbf{Forced opening.} A bid still sealed after the window is assigned to a solver, as in \cref{sec:bond}, drawn with the randomness of a later slot. A key matching $C$, from anyone, opens the bid, which is then checked as in a reveal and returns $B$; $d$ pays the solver. A failing tag, or a malformation claim that stands, excludes the bid and forfeits $d$ and $B$.
  \item \textbf{Settlement.} Once every bid is open or excluded, the contract sorts the valid bids by price, highest first, breaking ties by a hash of the sealed bid and the randomness of a later slot, fills them from $Q_k$ in that order, the last one partly, and charges every filled bid the lowest filled price. The rest of each escrow is refunded, and unsold units return to the seller. Later batches do not wait for this one.
\end{enumerate}

Settlement follows the reveal window at once if every bidder reveals. After a missed reveal, if the first solver assigned is honest and live and its moves are included in time, it waits from the assignment at most $D_s$ for an opening or a failing tag, $D_s + D_c$ for an unchallenged claim, and $D_s + D_c + 69\tau$ for a disputed one. With $D_s = D_c = 10$ minutes and $\tau = 2$ minutes, these are 10, 20 and 158 minutes. Each failed assignment can add as much again, and costs a staked solver its stake.
